% Natural Deduction Proof Template, revision 2.
%
% Build with:   pdflatex template-v2.tex && mv template-v2.pdf static/template.pdf
% The site serves it from static/ at /template; keep the two in step.
%
% Revised for legibility to the transcription pipeline (tools/transcribe.py and
% the /transcribe route in Web.hs). Every change answers a failure actually
% observed on the sixteen-page evaluation corpus, not a general principle.
%
%   * No shaded band above row (1). The grey header fill caused the most
%     persistent bug of the lot: the lone "1" in the first Depends cell sits
%     against that band and was repeatedly swallowed as part of the header, or
%     the whole first row was dropped. A double rule and a gap separate the
%     header instead.
%
%   * The Depends convention is printed in the column head. Blank dependency
%     cells were the one place the reader invented content, writing a "1" that
%     was not on the page. Naming the convention makes an empty cell a
%     deliberate mark rather than an apparent omission.
%
%   * A rule legend at the foot. Misread rule names were the commonest
%     remaining error -- "cP" for CP, "uI" for UI, "IMP" for MP. Printing the
%     vocabulary puts the canonical spellings in front of the student as they
%     write, and inside the photograph as the model reads.
%
%   * A wider Justification column. "1,2,3,4,5 vE" did not sit comfortably in
%     the old 2.5cm, and cramped writing is misread writing.
%
%   * Heavier rules, no grey anywhere. These are photographed on a desk in
%     ordinary room light; 0.3pt hairlines lose contrast badly.
%
%   * A version mark in the corner, so a photograph can be traced back to the
%     template that produced it as the evaluation corpus grows.

\documentclass[11pt]{article}
\usepackage[top=0.45in,bottom=0.4in,left=0.55in,right=0.55in]{geometry}
\usepackage{array}
\usepackage{amssymb}

% Row height: tall enough that a negation stroke or a subscript does not touch
% the rule above.
\renewcommand{\arraystretch}{1.82}

% Rules heavy enough to survive a phone photograph.
\setlength{\arrayrulewidth}{0.8pt}

\newcommand{\colhead}[1]{\textbf{\large #1}}

% If you would rather students marked "no dependencies" explicitly than left
% the cell blank, replace this with
%   \newcommand{\emptyhint}{\footnotesize\itshape write $\emptyset$ if none}
% An explicit mark is easier for a reader -- human or machine -- to tell from
% an unfilled cell. It departs from the convention in the book, so it is your
% call rather than mine.
\newcommand{\emptyhint}{\footnotesize\itshape (blank = none)}

\begin{document}
\thispagestyle{empty}

\noindent
{\Large\textbf{Natural Deduction Proof}}
\hfill
{\footnotesize Name:\ \rule{3.6cm}{0.4pt}\qquad Exercise:\ \rule{2.4cm}{0.4pt}}

\vspace{1mm}
\noindent\rule{\textwidth}{1.2pt}
\vspace{2mm}

% Outer rules as well as inner ones: a closed box makes the column
% boundaries unambiguous, which is what the reader keys on.
\noindent\begin{tabular}{|
  >{\raggedright\arraybackslash}p{2.6cm}|
  >{\centering\arraybackslash}p{1.1cm}|
  >{\raggedright\arraybackslash}p{9.5cm}|
  >{\raggedright\arraybackslash}p{3.5cm}|
}
\hline
\colhead{Depends} & \colhead{Line} & \colhead{Formula} & \colhead{Justification} \\[-3.2mm]
{\emptyhint} & & & \\
\hline\hline
& (1)  & & \\ \hline
& (2)  & & \\ \hline
& (3)  & & \\ \hline
& (4)  & & \\ \hline
& (5)  & & \\ \hline
& (6)  & & \\ \hline
& (7)  & & \\ \hline
& (8)  & & \\ \hline
& (9)  & & \\ \hline
& (10) & & \\ \hline
& (11) & & \\ \hline
& (12) & & \\ \hline
& (13) & & \\ \hline
& (14) & & \\ \hline
& (15) & & \\ \hline
& (16) & & \\ \hline
& (17) & & \\ \hline
& (18) & & \\ \hline
& (19) & & \\ \hline
\end{tabular}

\vspace{3mm}

\noindent\rule{\textwidth}{0.8pt}
\smallskip

\noindent{\footnotesize\textbf{Rules.}\ Write the rule exactly as printed here.
\ ($m$, $n$ stand for line numbers.)}

\smallskip
\noindent{\footnotesize
\begin{tabular}{@{}l@{\hspace{7mm}}l@{\hspace{7mm}}l@{\hspace{7mm}}l@{\hspace{7mm}}l@{}}
A                & $m,n$ MP  & $m,n$ MT   & $m$ DN            & $m,n$ CP \\
$m,n$ $\wedge$I  & $m$ $\wedge$E & $m$ $\vee$I & $m,n$ RAA     & $m$ QN \\
$m$ $\forall$E   & $m$ $\forall$I & $m$ $\exists$I & $m,a,n$ $\exists$E & $d,a_1,c_1,a_2,c_2$ $\vee$E \\
$m,n$ $\leftrightarrow$I & $m,n$ $\leftrightarrow$E & =I & $m,n$ =E & LEM \\
\end{tabular}}

\hfill{\tiny HLW-2}

\end{document}
